\documentclass[12pt,a4paper,oneside]{report}

\usepackage[utf8]{inputenc}
\usepackage[margin=1in]{geometry}
\usepackage{amsmath,amssymb,amsfonts}
\usepackage{graphicx}
\usepackage{booktabs}
\usepackage{hyperref}
\usepackage{cleveref}
\usepackage[numbers,sort&compress]{natbib}

\title{\textbf{Multi-Class Histopathology Semantic Segmentation on Whole Slide Images: From Supervised Baselines to Dynamic Attention Rectification}}
\author{\textbf{Pragyan Prayas Jena}}
\date{November 2026}

\begin{document}

\maketitle

\begin{abstract}
Whole slide imaging (WSI) in computational digital pathology produces gigapixel-scale tissue preparations requiring multi-class pixel-level semantic characterization. This thesis investigates multi-class tissue segmentation on the Breast Cancer Semantic Segmentation (BCSS) benchmark across tumor epithelium, stroma, necrosis, and lymphocyte infiltration. We establish a fully supervised U-Net baseline before exploring single-stage feature rectification and dynamic attention bias mechanisms.
\end{abstract}

\tableofcontents

\chapter{Introduction}
Introduction chapter placeholder.


\chapter{Literature Review}
Literature review placeholder.


\chapter{Methodology}
Methodology placeholder.


\chapter{Experiments and Results}
Experiments and results placeholder.


\chapter{Discussion and Conclusion}
Discussion and conclusion placeholder.


\bibliographystyle{plainnat}
\bibliography{references}

\end{document}
